\subsection{معماری و پیاده‌سازی}
پیاده‌سازی این ماژول در فایل \lr{Codes/framework/compression\_engine/proxsparse.py} قرار داشته و با شناسه‌های \lr{proxsparse} و \lr{proxsparse\_pruner} در رجیستری فشرده‌سازی ثبت گردیده است. ساختار معماری شامل موارد زیر است:
\begin{itemize}
    \item \textbf{کلاس ماژول \lr{ProxSparsePrunedLoRALinear}:} لایه هرس‌شده به عنوان جایگزین لایه‌های \lr{nn.Linear} و \lr{Conv1D} مستقر می‌شود و شامل وزن‌های منجمد، ماسک ماندگار (\lr{persistent buffer}) و شاخه‌های تنظیم‌پذیر ریزتنظیم \lr{LoRA} است.
    \item \textbf{پشتیبانی از ساختارهای سخت‌افزاری متنوع:} قابلیت پیکربندی برای هرس $2:4$ نیمه‌ساختاریافته سازگار با هسته‌های تنک پردازنده‌های امپر و هاپر (\lr{Ampere/Hopper Sparse Tensor Cores}) و هرس بلوکی با ابعاد سفارشی مانند $(4, 4)$.
    \item \textbf{مسیر خروج اضطراری (\lr{Fallback Mechanism}):} کنترل خطای بازسازی نسبی فروبنیوس پس از اعمال ماسک نهایی صورت می‌پذیرد؛ در صورتی که خطا از سقف آستانه \lr{max\_error} عبور کند یا اعداد نامعتبر پدید آید، وزن اصلی دست‌نخورده حفظ می‌گردد.
    \item \textbf{جداسازی لایه‌های بحرانی:} طبقه‌بندها، توکن‌های طبقه‌بندی، سرهای رگرسیون و بخش‌های کلیدی مدل از پروسه هرس مستثنی می‌شوند تا مدل دچار فراموشی فاجعه‌بار نگردد.
\end{itemize}
